% Guide paragraphs are added when the relevant section has been visually collated.
% Cross-references to later sections will be extended as their review progresses.
\paragraph{\S67. Aspheric maps replace contractibility as the organizing notion (pp.\ 188--191).}
Grothendieck reorients the modelizer program: the crucial objects are aspheric objects and weak equivalences, and its natural comparison maps are aspheric functors. He explicitly judges the earlier restriction to contractibility and special test functors to have been an unnecessarily long detour (p.\ 188), while retaining its possible relevance for ``canonical modelizers.'' The discussion now allows arbitrary small categories, rather than only test categories. For a basic localizer $\underline W$ he introduces $(\mathrm{Hot}^{W}_{A})$, also written $(\mathrm{Hot}_{A})$, as the localization of the presheaf category by $\underline W_A$; the construction is then organized by inverse image $i^*$ and the comparison morphism $\lambda_i(F):A_{/i^*F}\longrightarrow B_{/F}$, interpreted through the corresponding induced topoi (pp.\ 189--190). This prepares the promised description of asphericity structures and testing functors at the end of the section (p.\ 191); the next section interrupts that development. The dated marginal correction of 3 May on p.\ 189 expressly retracts the claim that Quillen's weak equivalences are precisely $\underline W_A$, saying that Quillen's are stronger. It is therefore essential to read the original assertion together with the margin. The witness changes handwritten ``or'' into equality in (1), drops the subscript $A$ on $\underline W_A$, omits the handwritten domain and codomain after the action on $(a,p)$ on p.\ 190, and prints the struck word ``real'' on p.\ 188. These differences concern notation, functorial context, or the author's revision; none licenses an unrecorded mathematical repair of the typescript.

\paragraph{\S68. Mathematical work, the Teichmüller tower, and the necessity of groupoids (pp.\ 191--195).}
Dated 27 May, the section explains the six-week interruption in the asphericity discussion through a retrospective of the author's changing mathematical interests. Its mathematical center is the ``Lego-Teichmüller construction game'': cutting and gluing surfaces, the tower of groupoids $T_{g,\nu}$, and the arithmetic action on their profinite completions (pp.\ 191--192, 194--195). The account of \emph{La Longue Marche}, tame topology and ``unfolding'' of stratified structures supplies the context for the anabelian project; the intervening reflection on creative interest, ambition and alternating periods of work belongs to the argument's historical presentation and must not be omitted as an aside (p.\ 193). The final page gives the technical lesson: reducing several natural base points to one can destroy visible symmetries and obscure a presentation by generators and relations, whereas groupoids preserve the geometry. Van Kampen situations supply the comparison, and the modular ``multiplicities'' are identified as topoi whose points have finite automorphism groups (p.\ 195). This is a return to the text's groupoid motivation, rather than a proof of the asphericity claims left pending in \S67; \S69 resumes a retrospective of the technical scratchwork. The section explicitly revises an earlier preference for groups or ``exterior groups'' in favor of groupoids (p.\ 195). The witness has several ordinary wording errors here---including ``case'' for ``care'' on p.\ 195---while the underlining and handwritten exclamation in ``topoi!'' preserve the emphasis of the typescript.

\paragraph{\S69. Derivators and integration of homotopy types (pp.\ 196--202).}
The technical retrospective explains why an ordinary derived or homotopy category retains too little structure: traces, determinants and the compatibility of iterated constructions motivate the system of localized diagram categories $D_A(I)$, contravariant in the indexing category $I$, with inverse-image functors expected to admit both adjoints (pp.\ 196--198). This develops the modelizer program recalled in \S67 and connects it with the geometric motivations of \S68. For categories, a covariant diagram can be integrated by forming its associated cofibered category; the comparison with localized slice categories yields the model for ``integration of homotopy types'' (pp.\ 198--199). The name ``derivator'' is introduced on p.\ 200, together with the internal stability expressed by $D(I\times J)$ and the aspiration to a corresponding homotopy theory of derivators. The geometric example is the system of tubes around strata of Mumford--Deligne compactifications: their $K(\pi,1)$ homotopy types, with Teichmüller-type groups, suggest an algebraic description from a diagram of fundamental groupoids (p.\ 201). A diagram in $(\mathrm{Hot})$ alone is explicitly said to be insufficient; a diagram in a model category is required. Grothendieck retracts his recent expectation that finite limits exist in $(\mathrm{Hot})$ (p.\ 201), and \S70 immediately examines the extra factorization condition proposed for the two adjoints (pp.\ 200, 202). Significant witness differences include $T$ replaced by $\mathscr D$ in the statement about inverse limits (p.\ 198), $q=f\circ p$ replaced by an equation beginning with $g$, and ``basic modelizer'' replaced by ``basic localizer'' (p.\ 199). These readings are clear on the images and are preserved rather than repaired from the witness.

\paragraph{\S70. Base change, duality, and the limitations of $(\mathrm{Cat})$ (pp.\ 202--206).}
Starting from the desired two-stage factorization $X\to Z\to Y$, Grothendieck asks for preservation of weak equivalences by base change along the second map and co-base change along the first, with either factor optionally a weak equivalence (p.\ 202). The goal supplies the extra condition proposed for derivator operations in \S69. The investigation of the base-change side leads to cohomological smoothness and properness of functors, characterized through asphericity; fibered categories are smooth, cofibered categories proper, and passage to opposite categories interchanges the two conditions (p.\ 203). The pairing of presheaves and copresheaves suggests a broader duality of topoi, but the claim that cosheaves on an arbitrary topos form a topos remains a question (pp.\ 203--204). The co-base-change side is less tractable: open or closed immersions in $(\mathrm{Cat})$ fail the hoped-for general assertion, whereas pushouts along monomorphisms in $A^{\wedge}$ have the required homotopy meaning through the integral of the associated span (pp.\ 204--205). The stated smooth/proper exceptions and the possible special case of an immersion that is also a weak equivalence must be distinguished from that rejected general assertion. This forces a return to elementary modelizers and Quillen factorization questions; the author finally accepts the Kan condition after describing his resistance to it, while retaining dissatisfaction with the reduction to semisimplicial methods (p.\ 206). These questions become part B of the program in \S71. The witness omits the monomorphism hook in the p.\ 204 square and replaces the actual objects $Y,X,Y'$ of the p.\ 205 integral diagram by categories of elements. The scan-confirmed diagrams preserve the distinction between these two levels of the construction.

\paragraph{\S71. A provisional program: modelizers, derivators, and topoi (pp.\ 206--212).}
The program dated 30 May organizes the preceding scratchwork into four tasks. Part A returns to the asphericity structures and canonical modelizers suspended in \S67; part B develops the base-change, co-base-change and factorization questions of \S70, with homotopy integrals, amalgamated sums and Mayer--Vietoris sequences as the connecting constructions (pp.\ 206--207). A handwritten transposition on p.\ 207 makes strong saturation one of the properties expected to follow from the first three basic-localizer axioms. Part C develops the derivator proposal of \S69 while asking what additional structure is needed for external Hom's and basic invariants. It then distinguishes group objects modeling pointed connected homotopy types from the truncated output of group objects in $(\mathrm{Cat})$, and asks which modelizers satisfy a Dold--Puppe theorem (pp.\ 207--209). The proposed trilogy consists of semisimplicial, cubical and hemispherical complexes, the last also called globular complexes by Brown; their boundary and degeneracy operations are related to regular cellular subdivisions of spheres. Multicomplexes give multiple chain complexes, so the simple Dold--Puppe statement does not extend unchanged. The classification of the three series and the need for enough degeneracy maps are presented as questions to investigate, not proved theorems (pp.\ 209--210). Part D, ``Back to topoi,'' seeks relative homotopy categories $D(T)$ and their possible opposite counterparts, extending the $f^*,f_!,f_*$ formalism from indexed categories to topoi. The discussion keeps existence of derived adjoints tentative and moves from objects to proobjects when a topos fails to determine an essentially constant homotopy type (pp.\ 210--211). Local asphericity then reconnects the program with \S35; the revised distinction between modelizing and strictly modelizing topoi is recalled from p.\ 68, and the section ends with the Lawvere-element condition for a finitely triangulated contractible space (p.\ 212). The next section begins the promised return to the asphericity game. The witness changes $f_*$ to $f^*$ and $Lf_!$ to a subscript-$1$ expression on p.\ 210, and applies a handwritten transposition there without marking it. The duplicated target-only display on p.\ 211 was an error of the draft transcription, absent from the scan; its removal restores the actual three-display sequence.

\paragraph{\S72. Criteria for aspheric functors and the first stability results (pp.\ 213--219).}
Returning to \S67 after the program of \S71, Grothendieck isolates the localizer assumptions needed for asphericity: mild saturation, the condition on categories with final object, and the locality condition Loc~3 (pp.\ 213--214). Proposition~1 identifies the comparison maps on categories of elements, their restrictions over objects, and the corresponding slice criteria; it then distinguishes these equivalent asphericity conditions from preservation of weak equivalences and equivalence after localization (pp.\ 214--216). The extra hypotheses in the implication diagram matter: the discrete-target example separates aspheric fibers from merely nonempty fibers, and the product-projection example shows that inducing the localized equivalence need not determine an aspheric map without the stated assumptions (p.\ 216). Under asphericity, the localized triangle commutes up to the canonical comparison isomorphism; this reconnects the discussion with the earlier ``inspiring assumption'' of \S28 without using that assumption as a proof (pp.\ 216--217). Proposition~2 transfers total asphericity along an aspheric functor and relates the inverse image of a separating homotopy interval to asphericity of the functor. Its handwritten addition makes the source category aspheric, and the long interlinear insertion specifies the homotopy condition (TH1) of p.\ 50; the corollary deduces strict test categories from a totally aspheric source and a local test target (pp.\ 217--218). Proposition~3 supplies composition and right cancellation, base change along a fibering map, and products (pp.\ 218--219). The base-change proof uses the fibration condition L5, whose derivation from the stated localizer assumptions is expressly postponed; Grothendieck explains that only part~(a) is needed for the immediate asphericity-structure formalism, preventing the announced application from relying on that deferred argument (p.\ 219). Remark~4 denies the opposite cancellation assertion and leads directly to the examples in \S73. The visual collation restores the proof-diagram arrow from (i'') to (iii), rather than to (i'), on p.\ 215; it also identifies the apparent deletion of the t in Hot on p.\ 214 as pen retracing. On p.\ 218 the comparison sign is a handwritten tilde over a bar, with no arrowhead, contrary to the draft transcription and its diff description. These are distinctions between what the image shows and what a mathematically plausible reconstruction might supply.

\paragraph{\S73. Adjoints, coasphericity, and cofinality (pp.\ 220--221).}
The continuation dated 11 June begins by suppressing the qualifying $\underline W$ when one basic localizer is understood. It sharpens Remark~4 of \S72: a functor from the one-point category into an aspheric category is aspheric exactly when its chosen object is initial, illustrating how much stronger a map condition is than asphericity of its source and target (p.\ 220). Proposition~4 enlarges the preceding isomorphism example: every equivalence is aspheric, and more generally a left adjoint is aspheric because its slice categories identify with slices having final objects. Passage to opposite categories then motivates the name ``coaspheric'' for the dual condition, satisfied by the right adjoint. The corresponding undercategories and their initial objects connect coasphericity with a promised all-dimensional cofinality statement, whose dimension-zero consequence is invariance of direct limits (pp.\ 220--221). The analogy from \S70 returns: asphericity resembles properness with aspheric fibers, and coasphericity resembles smoothness with aspheric fibers. The comparison is explicitly weakened because arbitrary base change does not preserve either condition; the applicable base changes are fibering/smooth or cofibering/proper, respectively. Even an equivalence of categories can have an empty fiber, so its simultaneous asphericity and coasphericity do not supply arbitrary base-change stability. Universal weak equivalences, called ``trivial Serre fibrations'' here, are proposed as a stronger common condition and assigned to a later systematic investigation (p.\ 221). The section then turns to the overdue product discussion of \S74. On p.\ 220 the crossed handwritten heading above Indeed is Remark~5, not the draft's tentative Proof; it is written again at the subsequent paragraph. Whether the placement strokes also delete its typed initial The remains unresolved and is marked as such. The p.\ 221 witness changes singular condition to plural conditions and supplies a title for \S74 absent from the typescript.

\paragraph{\S74. Products, diagonal functors, and canonical comparison of models (pp.\ 221--230).}
The overdue product discussion supplies concrete constructions for the abstract stability results of \S72. The external product of presheaves has a category of elements identified with the product of the two categories of elements; it preserves aspheric objects and weak equivalences in the stated variables, and yields the product assertions for locally aspheric and totally aspheric categories (pp.\ 221--222). The corollaries distinguish local test, test, and strict test categories: the hypothesis on the other factor changes from arbitrary to aspheric to totally aspheric. Mixed multicomplexes illustrate the construction (p.\ 223), reconnecting it with the trilogy in \S71. The comparison functor associated with an aspheric map between modelizing categories leads to a functor from $\mathrm{Asph}(A,B)$ into a category of comparison functors and factors through its fundamental groupoid. Nonemptiness and connectedness suffice for the existence of an isomorphism class of comparison, while a canonical transitive system requires the stronger simple-connectedness conclusion; the text carefully separates these claims (pp.\ 224--225). The diagrams on p.\ 226 distinguish categories, presheaf modelizers, and localized homotopy categories, including the middle presheaf triangle omitted from the witness.
The diagonal criterion then identifies total asphericity of $A$ with asphericity of $A\to A\times A$, and permits simultaneous comparison through products of strict test categories (p.\ 227). Proposition~4 treats pointwise products of functors into a category $B$ with binary products and aspheric presheaves $\underline{\mathrm{Hom}}(b,y)$: the total-asphericity assumption on $A$, moved by a handwritten transposition into the hypotheses, is essential to the stated construction (p.\ 228). The external-product corollary is sharpened by revision to an equivalence: when $i$ is aspheric, $i\boxtimes i'$ is aspheric if and only if $A'$ is aspheric. Its proof uses a lemma on a presheaf over $A\times A'$ whose restrictions to $A$ are aspheric. The resulting category-of-elements projection is a weak equivalence by the fibration condition L5, explicitly recalling the claim and deferred justification on p.\ 219; this dependency should remain visible (p.\ 229). Contractibility of the objects of $B$ is discussed as a stronger practical hypothesis. Finally, closure of $\mathrm{Asph}(A,B)$ under binary products makes its nonempty case connected with the needed homotopy strength, giving the canonical transitive system of isomorphisms between the localized functors and the canonical functor announced on p.\ 225 (p.\ 230). A handwritten margin points forward to the generalization in \S82, p.\ 272. The image-based edition retains the unusual $X\longrightarrow B_{/X}$ display on p.\ 224, the missing slash in the final category-of-elements notation on p.\ 226, and the visibly typed product x and lowercase subscript o on p.\ 228. The hemisphere-like symbol on pp.\ 223 and 229 remains a doubtful approximate glyph; neither mathematical expectation nor the witness settles its exact drawn form.

\paragraph{\S75. Asphericity structures, testing functors, and weak equivalences (pp.\ 230--237).}
The generalities of \S\S72--74 now become a structure on a possibly large $\underline U$-category $M$. After making the universe convention explicit, Grothendieck defines a class $M_{as}$ of aspheric objects through a small category $A$ and a functor $i:A\to M$ that lands in $M_{as}$ and detects it by the presheaf $i^*(x)=\mathrm{Hom}(i(-),x)$ (pp.\ 230--231). Such an $i$ is a “testing functor.” Every small full subcategory supplies an asphericity structure, and Proposition~1 plus its corollary shows that any structure can be generated this way; enlarging the generating full subcategory within $M_{as}$ preserves the structure (p.\ 232). The following terminology discussion distinguishes testing functors from the earlier test functors of pp.\ 175--176. The essential condition depends on the possibly large category $M_{as}$ itself, rather than on its embedding in $M$, so the author tentatively adopts $M_{as}$-aspheric, or simply aspheric, as a synonym. His use of both “aspherical” and “aspheric” in this discussion is preserved (p.\ 233).
Proposition~2 defines $W=(i^*)^{-1}(\underline W_A)$ and proves independence of the testing functor by passing to a common small full subcategory. Mild saturation is immediate; strong saturation remains conditional on the corresponding result for the basic localizer (p.\ 234). When $M$ has a final object, calling the structure aspheric means that this final object belongs to $M_{as}$; the resulting criterion recovers aspheric objects from the maps to the final object. The small-generator descriptions then characterize both the admissible weak-equivalence sets and the asphericity classes. Without the aspheric final-object hypothesis, a weak-equivalence set need not determine a unique asphericity structure. The analogy with contractibility and homotopism structures in \S\S51--52 motivates a promised construction from contractibility to asphericity and a uniqueness assertion for the intended canonical examples (p.\ 235). Proposition~4 recasts the criteria of p.\ 214 as six conditions: three about aspheric objects and three about weak equivalences. The object conditions are equivalent and imply the others; the revised converse requires both $A$ and the asphericity structure to be aspheric. All six conditions become equivalent when the specified generating category $B$ contains the final object. If no $B$ is fixed, enlarging one to contain $i(A)$ and the final object yields equivalence of the first five conditions (pp.\ 236--237). These added hypotheses, some typed interlinearly and some handwritten, are central to the proposition and were partly concealed by malformed insertion macros in the draft. Section~76 immediately tests the formalism on examples.
Two visible mathematical slips are retained for the reader to assess: the universe paragraph says “two objects of $U$” where the witness substitutes $M$ (p.\ 230), and Proposition~2 says “Then $M$ is” a mildly saturated subset of $\mathrm{Fl}(M)$ after defining $W$, where the witness prints $W$ (p.\ 234). Proposition~3 and its corollary also call the relevant map “aspheric,” while the explicit revisions of Proposition~4 replace that word by “a weak equivalence” (pp.\ 234, 236); the edition does not silently harmonize the earlier wording. Several correction layers remain uncertain, including the usual-localizer subscript at the preceding section boundary and overtyped letters in these definitions and proofs.

\paragraph{\S76. Basic examples and total asphericity (pp.\ 237--239).}
Three examples test the definitions of \S75. For $M=(\mathrm{Cat})$, the $\underline W$-aspheric categories form the basic asphericity structure, generated through $a\mapsto A_{/a}$ for a weak test category $A$. The shift from the basic modelizer to this structure puts aspheric objects before weak-equivalence arrows. The distinction matters for functors into $(\mathrm{Cat})$: asphericity supplies the induced localized functor, whereas the weak-test requirement asks for a homotopy-model comparison. Theorem~1 of p.\ 176 suggests equivalence of these conditions when the values are contractible, and the author identifies a comprehensive answer to the “silly question” of \S46, p.\ 95, as the motivation for the present formalism (pp.\ 237--238). The second example takes topological spaces with the weakly contractible spaces as aspheric objects; the geometric realization of simplices supplies the testing functor, and the text appeals to Gabriel--Zisman for the stronger modelizer statement. Grothendieck adds the previously omitted assertion that the first two structures are aspheric and later notes that the resulting weak equivalences are the usual ones (p.\ 238). The third example is $A^\wedge$ with its usual aspheric presheaves, generated by the canonical fully faithful inclusion of any small $A$. Its weak equivalences are $\underline W_A$, while its asphericity is equivalent to asphericity of $A$; the modelizing property is not asserted for arbitrary $A$. A functor $A'\to A$ is aspheric precisely when its composite into this presheaf structure is $M_{as}$-aspheric (pp.\ 238--239). This leads to “totally aspheric” structures: $M$ admits finite products and finite products of aspheric objects, including the empty product, are aspheric. Proposition~5 identifies this condition with the existence of a small generating subcategory stable under finite products, and with a testing functor from a category whose presheaf category is totally aspheric. It reconnects the product arguments of \S74 with the general formalism and prepares the localized canonical functor of \S77 (p.\ 239). The visual check restores a missing typed a before the repaired l in “realization,” retains the literal thatnthis and appheric typing, and distinguishes typewritten interlinear additions from pen text. The upward margin arrow at the end of the examples had been misrepresented by three brackets in the draft; its direction is restored, without attributing a mathematical meaning to the layout mark.


\paragraph{\S77. The canonical homotopy functor and the meaning of modelizing (pp.\ 239--241).}
An asphericity structure now determines more than its weak-equivalence class: it supplies a canonical functor from $\mathrm{Hot}_M=W^{-1}M$ to $(\mathrm{Hot})$. Choosing a testing functor $i:A\to M$, Grothendieck uses the composite through $A^\wedge$ and its category-of-elements functor; the inverse-image descriptions of the weak equivalences permit localization, while conservativity is stated conditionally on strong saturation (pp.\ 239--240). To compare two choices, he places their images in one small full subcategory $B$ of $M_{as}$, obtains the comparison isomorphisms from aspheric functors, and then enlarges $B$ to prove independence. A common $B$ for three choices gives transitivity, completing the canonical construction promised in \S74. Formula (13) on p.\ 240 actually names its composite $i_M$, while the following localized functor carries the extra choice index; the witness's harmonization to $i_{M,i}$ in (13) is not adopted. The two sloping comparison signs there are hand-drawn bars with tildes, without arrowheads. The new definition calls the structure modelizing when this particular canonical functor is an equivalence; this is stronger than merely requiring an abstract equivalence between $W^{-1}M$ and $(\mathrm{Hot})$. In the presheaf example it says exactly that $A$ is a pseudo-test category (p.\ 241). The author immediately hesitates over the terminology and proposes a parallel to the hierarchy of pseudo-test, weak test, test, strict test categories and contractors, but leaves the final qualification as question marks. His deletion of “progressively stronger” in that discussion remains visible; the witness prints the deleted words as live text. Section~78 therefore begins with the modelizing name expressly provisional, and uses the canonical construction to revisit the observation and silly question of \S46 rather than treating the proposed hierarchy as already settled.


\paragraph{\S78. Test functors separated into asphericity and modelizing conditions (pp.\ 241--246).}
The provisional definition of modelizing from \S77 supplies the setting for a reformulation of the observation and silly question of \S46. Theorem~1 assumes a modelizing structure $(M,M_{as})$, a pseudo-test category $A$, a functor $i:A\to M$ with aspheric values, and the existence of a final object of $M$; the last hypothesis is a typed interlinear addition, not an original underlined parenthetical sentence. It equates the asphericity and weak-equivalence criteria of Proposition~4, p.\ 236, with the condition that $i^*$ induce an equivalence between the localized categories (p.\ 242). A lemma says that a final object in a modelizing structure is aspheric. Its supporting corollary characterizes $M_{as}$ as the objects whose images under the canonical map to $(\mathrm{Hot})$ are final. The proof uses strong saturation of the basic localizer, announced as a consequence of Loc~1--3 but deferred to part~V; the author explicitly offers the additional asphericity assumptions on the two final objects as a way of avoiding the feared circularity before that proof is supplied (pp.\ 242--243). The canonical triangle from \S77 then proves the localized equivalence. The proof's literal reference to conditions through (ii') is retained even though the statement speaks of six conditions through (ii'').
The theorem prompts a revision of the terminology inherited from \S65. The evolving pseudo-test-functor definition first describes model-preserving composites, then separates the intended conditions into: $A$ is pseudo-test, $(M,M_{as})$ is modelizing, and $i$ is $M_{as}$-aspheric. The discussion revises both its hypotheses and its formulation on the page; these are stages of the argument rather than interchangeable completed definitions (pp.\ 243--244). Passing to weak test functors strengthens the condition on $A$ to weak test. Passing to test functors requires a test category and asphericity both of $i$ and of all its slice-induced functors. Definition~2 calls the latter family of slice conditions locally aspheric, and the conjunction with asphericity of $i$ totally aspheric. The local criterion additionally requires that $i$ factor through $M_{as}$ and that $i^*(x)$ be aspheric over the final presheaf, with $x$ restricted to a generating subcategory containing $i(A)$ when desired (p.\ 245). The Remarks paragraph draws the conceptual conclusion: the independent asphericity conditions have wider relevance than the combined test-functor terminology. Yet the earlier theorem on $(\mathrm{Cat})$ still has one advantage, namely its criterion involving only $i^*(\Delta_1)$ under the contractibility hypothesis on the values of $i$. Recovering such a minimal criterion requires relating asphericity structures to the contractibility structures of \S51, which becomes the subject of \S79 (p.\ 246). The typed minimal is inserted before criterion without deleting generating in the preceding line. Most of the final sentence on p.\ 246 is physically covered by a laid-over sheet; its tentative transcription remains doubtful and is not treated as fully recovered evidence.


\paragraph{\S79. From contractibility to asphericity and the general test-functor criterion (pp.\ 247--255).}
The entry dated 20 June supplies the missing bridge identified at the end of \S78. Proposition~1 removes a smallness assumption on a generating full subcategory $N$: the existence of a small generating part, a small inclusion detecting asphericity on $N$, and the assertion that $(N,N)$ is an asphericity structure are equivalent. The name “aspherator” emphasizes that the condition belongs to $N$ itself, while its dependence on the chosen universe and basic localizer remains explicit (p.\ 247). The author then strengthens the standing convention for contractibility structures to require a small generating full subcategory and recalls the finite-product assumption on $M$. Proposition~2 makes $M_c$ generate a totally aspheric structure $M_{as}$; its clause about a chosen generating category $C$ acquires the additional handwritten hypothesis that $C^\wedge$ be totally aspheric (p.\ 248). The proof uses finite-product closure, contractibility of the induced presheaves, and Proposition~5 of p.\ 239. The following remark distinguishes this construction from arbitrary asphericity structures and asks about a canonical largest contractibility structure generating a given one. Its claim about every previously encountered modelizing structure is immediately retracted using nonstrict test categories, then narrowed to modelizing totally aspheric structures (pp.\ 249--250). The second part recalls locally and totally aspheric presheaves and the homotopy interval structure generated by the locally aspheric ones. The source's reference to part IV, \S\S54--55, is retained as written. Six equivalent compatibility conditions H1--H4' express preservation of homotopies, homotopisms, intervals, or contractible generators by a product-preserving functor; Remark~2 explains exactly what generation by $C$ means (pp.\ 250--252). Theorem~1 applies this to $i^*:M\to A^\wedge$ when $i:A\to M$ lands in $M_c$: compatibility with the homotopy structures, local $M_{as}$-asphericity, and the same presheaf criterion restricted to $C$ are equivalent. Its corollaries add asphericity of $A$ for total asphericity, and the test-category and modelizing requirements for a test functor. A separating multiplicative interval supplies a sufficient criterion forcing $A$ to be a local test category; later margin notes question smallness of $C$ and say that multiplicativity is unnecessary, pointing to p.\ 348, Remark~1 (pp.\ 252--253). The proof abandons a boxed localization argument in favor of a lemma and corollary about testing on generators; the blank relation positions in the deleted draft are not completed editorially (p.\ 254). Remark~3 specializes to totally aspheric $A$ and obtains a strict test category under the interval hypothesis, closing the promised restatement of the earlier result. Its corollary reference is corrected by hand from 2 to 3, and a separate margin note points to a stronger version on p.\ 347 (p.\ 255). Section~80 then reverses the direction of inquiry, asking what homotopy structure is recovered from the associated weak equivalences.



\paragraph{\S80. Recovering homotopy from weak equivalences and the canonical-modelizer proposal (pp.\ 255--258).}
The construction $M_c\mapsto M_{as}$ from \S79 is set beside the reverse-looking association $W_a\mapsto h_{W_a}$ of \S54. The first direction distinguishes contractibility from the weaker homotopy interval and homotopism structures; it also makes clear why the associated weak equivalences $W_a$ should not be confused with the generally smaller class $W_c$ of homotopisms. The author recalls how the aspheric-object viewpoint emerged after the earlier “awkward approach,” and corrects and supplements the references to that development (p.\ 255). For an asphericity structure with aspheric final object, $W_a$ and $M_{as}$ determine each other. The intervals generating $h_{W_a}$ are those for which every projection $x\times I\to x$ lies in $W_a$. Under total asphericity this is equivalent to asphericity of $I$, an observation promoted from a tentative expectation to an immediate statement. Two questions then ask whether passing from $M_c$ to $M_{as}$ and back recovers the same homotopy structure, and which extra conditions on $M_{as}$ ensure that the recovered structure comes from a contractibility structure generating it (p.\ 256). A third question, now redundant with the preceding observation, is boxed and deleted; its formula reads $UW_a$, with both letters visible (p.\ 257). Before answering, Grothendieck resumes the “idyllic picture” of canonical modelizers from \S50, incorporating \S\S57 and 59. A finite-product category $M$ with a functor $\pi_o:M\to(\mathrm{Sets})$ receives an interval structure from projections inducing bijections under $\pi_o$. If this structure comes from $M_c$, then $M_{as}$, $W_a$, and the canonical functor $W_a^{-1}M\to(\mathrm{Hot})$ follow. The further compatibility between $\pi_o$ and the composite through $(\mathrm{Hot})$ remains to be imposed; the mild localizer condition L~a) on p.\ 165 is needed to define the final connected-components functor (pp.\ 257--258). When $M$ is totally 0-connected, $\pi_o$ is its intrinsic connected-components functor, and the resulting asphericity structure is modelizing, the construction deserves the name canonical modelizer because it depends on the category itself. The closing Remark calls the definition essentially complete while explicitly retaining the compatibility question. This unresolved point, and the two recovery questions, lead into \S81. The visual collation distinguishes the typed interlinear hypotheses from their pen locators and preserves the literal “homotopisms into homotopisms” on p.\ 257, where the witness substitutes “isomorphisms.”




\paragraph{\S81. Recovering contractibility and comparing the three families of notions (pp.\ 259--266).}
The entry dated 25 June announces the positive answer to the first recovery question of \S80: a contractibility structure $M_c$ on a category with finite products is recovered from its associated asphericity structure by taking the objects contractible for the homotopy interval structure generated by all aspheric intervals. Thus $\mathrm{Cont}(M)\to\underline{W}\text{-}\mathrm{Asph}(M)$ is injective when the basic localizer satisfies Loc~4; the proof is deferred to \S82. The four homotopy notions recalled in (2)--(4) place this assertion beside the earlier inclusions from contractibility to interval structures and homotopism structures, with the latter equivalent to homotopy relations (p.\ 259). Changing the basic localizer from $\underline{W}$ to a coarser $\underline{W}'$ induces a canonical change of asphericity structure, and the triangle on p.\ 260 makes this construction compatible with the inclusion of contractibility structures. Loc~4 requires weak equivalences of small categories to induce bijections on connected components. Among localizers satisfying it, $\underline{W}_o$ is the coarsest; the finest basic localizer $\underline{W}_{\mathrm{oo}}$ is introduced as their intersection and its identification with the usual weak equivalences is promised for part V. The strongest recovery statement is therefore the elementary case associated with $\underline{W}_o$ (pp.\ 260--261). Writing $M_o$ for the 0-connected objects and $W_o$ for the corresponding maps, (17)--(20) show how either determines the other and how both are recovered from the connected-components functor $\pi_o:M\to(\mathrm{Sets})$. The first line of (19) literally omits “is” before “a one-point set.” A handwritten marginal note objects to the name 0-connected maps and proposes 0-équivalences, but its crowded continuation remains partly unreadable; the later asterisk on p.\ 266 refers back to this objection (p.\ 262). The retrospective discussion explains how the special canonical connected-components functor delayed recognition of this general construction and why much of part III is superseded by part IV, while its homotopy notions remain essential (pp.\ 263--264). The three families are test notions, homotopy notions, and asphericity notions. Test and asphericity notions depend on the basic localizer; homotopy notions are absolute. Their order in a deductive account would reverse much of their order of discovery. Contractibility becomes the hinge between homotopy and asphericity: it is both the strongest of the four homotopy structures and a specially constrained asphericity structure (pp.\ 264--265). The revised displays then enumerate ten mutually determining data, including $M_c,J_c,W_c,R_c$, the 0-connected and $\mathrm{oo}$-aspheric objects and maps, and the two canonical functors to sets and homotopy types. This count is reached through visible revisions from four to three asphericity-flavored data and from eight through nine to ten structures. Two further relations $R_o,R_{\mathrm{oo}}$, and more generally $R_{\underline{W}}$, compare maps having equal images under the corresponding functors; whether such a relation alone recovers $M_c$ is explicitly left unclear. The inclusions and commutative comparison in (26), followed by the localization chain in (27), complete the section (pp.\ 265--266). The checked source preserves the typed “oo” notation, identifies the underlined $\pi_o$ in (26), and distinguishes pen arrows from typed diagram objects. Its revision record also retains the five-dot reinstatement of “being” on p.\ 259, the phrase transposition on p.\ 260, and the cancelled draft displays on p.\ 265.



\paragraph{\S82. Proof of recovery and consequences for internal Hom objects (pp.\ 267--272).}
Theorem~1 supplies the proof promised in \S81. Under Loc~4, with finite products and the stipulated small full generating category, the contractible objects for the homotopy interval structure generated by all aspheric intervals are exactly the original $M_c$. The argument compares $M_c$, the recovered contractible objects, and $M_a$, then reduces the essential point to Lemma~1: any two sections of an aspheric object are already $M_c$-homotopic (pp.\ 267--268). Passing to the elementary localizer and a small full generating category $C$ reduces the problem to its presheaf topos. Lemma~2 restates the finite-chain argument from \S60, p.\ 149: in a totally 0-connected topos with the stated section and generator assumptions, two sections of a 0-connected object can be joined through a finite chain of objects of $C$. The proof of this lemma is left to the reader. The displayed inclusion $\underline W_o\subset\underline W$ on p.\ 268 is retained as printed, although the witness reverses it; it should be read alongside the coarseness discussion on pp.\ 260--261.
The next application assumes the relevant internal Hom objects exist. For a pointed object $X$, the endomorphism object $\underline{\mathrm{Hom}}(X,X)$ has the identity and the constant map as endpoints and supplies a multiplicative interval. The resulting equivalences connect contractibility of $X$, contractibility or 0-connectedness of internal Hom objects, and homotopy equivalence of the constant-map inclusion $Y\to\underline{\mathrm{Hom}}(X,Y)$ (pp.\ 269--270). These are consequences of the recovered homotopy structure, not new definitions of the asphericity class. The source uses $\pi_o,\pi_1$ for the endpoints in one display; that notation is preserved. Lemma~3 compares homotopy classes of maps from $Y\in C$ with the corresponding classes after applying $i^*$, referring back to p.\ 144. Together with the availability of sections of contractible generators, this gives Proposition~1: homotopy classes of sections map bijectively to connected components. Its internal-Hom corollary identifies homotopy classes of maps $X\to Y$ with the connected components of $\underline{\mathrm{Hom}}(X,Y)$ when that object is representable (pp.\ 270--271).
Proposition~2 generalizes the product argument of \S74, p.\ 228, under internal-Hom hypotheses. The revisions add total asphericity of the indexing category and alter several occurrences of the generating category and object variables; the cancelled final criterion remains visible in the diplomatic account (p.\ 271). Setting $B=C=M_c$ yields the corollary that the product of an aspheric functor with a functor valued in $M_c$ is aspheric. The author explicitly wonders whether the internal-Hom existence assumption for pairs of contractible objects is superfluous. Hence the category of aspheric functors valued in $M_c$ is stable under binary products. If nonempty, it gives the canonical functor $\mathrm{Hot}_{M,\underline W}\to\mathrm{Hot}_{A,\underline W}$, defined up to unique isomorphism through the transitive system recalled from p.\ 230 (p.\ 272). This closes the recovery proof and its product application; the next section asks for a converse characterization of the asphericity structures to which it applies.

\paragraph{\S83. A tautological characterization and an unresolved functor question (pp.\ 272--273).}
Grothendieck returns to the question omitted from the preceding theorem: which asphericity structures arise from contractibility structures? He first imposes total asphericity, then forms the interval structure $J$ generated by intervals whose objects lie in $M_a$ and calls its contractible objects $M_c$, warning that this set need not itself be a contractibility structure (pp.\ 272--273). The two conditions are that $M_c$ generate $J$, equivalently that two sections of any object of $M_a$ be $M_c$-homotopic, and that $M_c$ generate the given asphericity structure through a small full subcategory with aspheric inclusion. The first condition ensures that $J$ does arise from a contractibility structure, necessarily this $M_c$. The characterization is expressly called a tautology. A parallel request asks when a functor $\phi:M\to\mathrm{Hot}_{\underline W}$ is isomorphic to one associated with a contractibility structure, especially when it reduces to $\pi_o:M\to(\mathrm{Sets})$. The text offers no useful necessary-and-sufficient criterion and leaves open whether suitable exactness and internal-Hom assumptions would produce one. Its final underlined question asks whether such a functor can have nontrivial automorphisms, connecting this inquiry to the “inspiring assumption” of \S28 (p.\ 273). The checked text preserves the visible overtyped ending of “generated,” the parentheses around Sets, and “writing it down”; it does not replace these with the witness's readings.


\paragraph{\S84. A provisional notion of morphism and the direct-image problem (pp.\ 274--280).}
The entry dated 29 June begins with letters to Faltings and Porter, then describes the unexpectedly rapid emergence of a proposed morphism notion after prolonged scratchwork (p.\ 274). For two asphericity structures and a functor $f^*:M\to M'$ with left adjoint $f_!$, the proposition parallels the six conditions of \S75, p.\ 236: preservation and detection of aspheric objects, and corresponding assertions about weak equivalences. Its extra assumption (Awk) requires an aspheric testing functor into $M'$ whose composite with $f_!$ lands in the chosen generating part of $M$. The page begins with a generating category $B$ but uses $C$ in that auxiliary clause; the source wording is retained. A triangle through $A^\wedge$ reduces the equivalences to the earlier proposition. The revised conclusion first equates five conditions under the aspheric-final-object hypotheses, then all six if the final object belongs to the chosen generating category (p.\ 275). The canonical localization functors give a second triangle, and the corollary proves that a morphism between modelizing structures induces an equivalence of localized categories. The author calls this the answer to the “silly question,” with the source reference \S45, p.\ 95 (p.\ 276).
The rest of the section examines the restrictions rather than treating the proposed definition as settled. The left-adjoint requirement is compared with preservation of inverse limits. Condition (Awk) is reformulated through a small generating subcategory of $M'_a$ taken by $f_!$ into $M_a$, or through the generating property of $(f_!)^{-1}(M_a)\cap M'_a$. Replacing it by the stronger condition $f_!(M'_a)\subset M_a$ would pair it neatly with $f^*(M_a)\subset M'_a$, but its reasonableness is unresolved (pp.\ 276--277). Question~1 asks whether the left adjoint $i_!:A^\wedge\to M$ attached to an aspheric functor preserves aspheric objects. In the presheaf case this becomes Question~1': when does $i_!:A^\wedge\to B^\wedge$ preserve them? Dropping asphericity of $i$ reduces the issue to the colimit of the representables $i(a)$ over an aspheric category. A cofibering example with 0-connected fibers shows why asphericity of the indexing category alone does not suffice; asphericity of $B$ must at least enter the discussion (pp.\ 277--278).
The proposed counterexample on p.\ 279 concerns an amalgamated sum $a\amalg_c b$. The monomorphism case recalls \S70, but the author explicitly has not completed a general counterexample. Neither the extra initial-object condition nor total asphericity appears convincing to him; imposing a strict test category leads to a further question he does not know how to handle. The checked edition preserves the inward arrow directions of diagram (9), despite the surrounding reference to a dual category, and the capital $B$ in the inline map $c\to B$. Finally, for $M=(\mathrm{Cat})$, the canonical slice functor $a\mapsto A_{/a}$ supplies the positive example when $A$ is a weak test category, with both associated functors modelizing and the left adjoint preserving aspheric objects. The standard inclusion $\Delta\hookrightarrow(\mathrm{Cat})$ is the decisive unresolved test: its left adjoint is not generally modelizing, and preservation of aspheric objects still has to be checked (p.\ 280). At the opening of \S85 on the same page, Grothendieck retracts the provisional morphism notion because stability under composition has not been established; this later objection governs how the preceding proposition and corollary should be read.

\paragraph{\S85. Product preservation and the limits of left exactness (pp.\ 280--285).}
The entry dated 30 June begins with Porter's reprints and the author's surprise at Thomason's announced closed model structure on categories. It immediately retracts the preceding proposed morphism notion: stability under composition has not been secured. Requiring $f_!$ to preserve aspheric objects would supply that stability, but the standard inclusion $\Delta\hookrightarrow(\mathrm{Cat})$ remains a difficult test. The author already knows that its left adjoint does not preserve all weak equivalences; he now asks whether it preserves aspheric or at least contractible objects (pp.\ 280--281). Finite-product preservation would prove the latter because the images of the generating simplices are contractible. Proposition~1 supplies the needed criterion for a cocomplete $M$ and $f:A\to M$: under commutation of the product functors with all small direct limits, the binary-product comparison for $f_!:A^\wedge\to M$ need only be checked on representables. The hypothesis is visibly strengthened from distribution over arbitrary sums. Under the corresponding base-change hypothesis, the fibered-product criterion reduces the two factors to representables, while retaining a general presheaf as base; preservation of a final object then gives the stated left-exactness corollary (p.\ 281).
The abandoned proof on p.\ 282 leaves unfinished relation positions inside its cancellation. The replacement proof writes presheaves as direct limits of representables and uses the displayed interchange of products with two direct limits. Passing to slice categories handles the fibered-product part. Remark~1 then identifies a limitation: reducing the base itself to an object of $A$ does not yet establish preservation of all fibered products without an additional monomorphism argument; a margin note points forward to \S89 (pp.\ 282--283). Proposition~2 assumes full faithfulness of $f$ and generation of $M$ by strict epimorphisms, expressed by $P\simeq\mathop{\mathrm{colim}}_{A_{/P}}f(x)$. It obtains finite products and the fibered products whose base comes from $A$. The proof uses the full-subcategory identification of the relevant slice categories. Remark~2 refers to SGA~4, volume~1, expos\'e~I, \S7, especially Proposition~7.2 on p.\ 47, for the connection between strict generation and full faithfulness of $f^*$ (pp.\ 283--284).
For $(\mathrm{Cat})$, Proposition~3 applies these exactness hypotheses and gives the finite-product conclusion together with the restricted fibered-product assertion. The nerve provides the familiar $\Delta$ example, but the extension to every full subcategory containing $\Delta$ carries the manuscript reservation “hasty!” (p.\ 284). Remark~3 doubts full left exactness even for $\Delta$, appealing to the generators-and-relations description cited to Gabriel--Zisman and the difficulty of preserving monomorphisms. The later margin note explicitly says that $i_!$ is not left exact (p.\ 285). Section~86 therefore proceeds from finite-product preservation and its contractibility consequence, rather than treating the stronger left-exactness claim as proved. The checked source retains literal objedt and the unlabelled pen curve beside hence; neither is silently replaced by the witness's smoother reading.



\paragraph{\S86. Bimorphisms of contractibility structures and the remaining morphism problem (pp.\ 285--290).}
Finite-product preservation from \S85 gives $i_!(A^\wedge_c)\subset(\mathrm{Cat})_c$ for the specified strictly generating subcategories of contractible categories containing $\Delta_1$. The text distinguishes this from the doubtful preservation of all aspheric objects. When both asphericity structures arise from contractibility structures, the condition $f_!(M'_c)\subset M_c$ suffices for the earlier purpose and is stable under composition. The results of \S81 and Loc~4 also allow it to be tested by the apparently weaker requirement that the relevant generators map to 0-connected objects (pp.\ 285--286). This replaces the troublesome (Awk) of \S84 and returns the six equivalent asphericity criteria. The resulting condition can be expressed through the induced functor between the aspherators of contractible objects (p.\ 286).
The theorem on p.\ 252 converts the corresponding inverse-image condition into $f^*(M_c)\subset M'_c$, producing a definition independent of the basic localizer. The author calls an adjoint pair satisfying these two inclusions a “bimorphism,” and changes the heading Proposition to the handwritten “Scholie” (p.\ 287). Part (a) reduces the inverse-image inclusion to testing generators for 0-connectedness; the analogous direct-image reduction additionally assumes finite-product preservation. Part (b) says that a bimorphism makes $f^*$ both preserve and detect the aspheric objects and weak equivalences for each basic localizer. Part (c) gives the triangle with the canonical functors of \S77, and hence an equivalence on localizations when both structures are modelizing. Part (d), assuming the direct-image inclusion and Loc~4, equates the bimorphism condition with the indicated inclusions, asphericity of $f_{!c}$, and preservation of an aspheric testing functor valued in $M'_c$ (p.\ 288). The checked page retains its literal $M_{\underline W'}$ in the opening notation and $\underline W_M$ in the final line. The triangle's handwritten $f^*$ has no bar, although the nearby inserted localized functor does.
The comments call the Scholie a long tautology in light of \S\S81--82, but emphasize the care required to separate these relationships. A product-preserving $f_!$ also acts on strict homotopy localizations; Grothendieck doubts the geometric relevance of that localized map except in the exceptional situation where both adjoints preserve the weak-equivalence models and induce inverse equivalences. The Scholie replaces the “seducing” proposition of \S84 as his answer to the question of \S46, almost 200 pages earlier. Whether finite-product preservation should enter the definition itself remains a terminological question (p.\ 289). More fundamentally, a satisfactory notion of morphism for arbitrary asphericity structures is still missing, as is a convincing way to dispense with the left adjoint. The closing paragraph explicitly shifts the emphasis toward contractibility structures and entertains the possibility that asphericity structures are principally an important by-product of them (p.\ 290). These reservations qualify the confident language about the bimorphism definition without cancelling it.



\paragraph{\S87. Thomason's Dwyer maps and the independence of homotopy limits (pp.\ 291--292).}
The entry dated 1 July reports an overnight reading of Thomason's preprint. Grothendieck welcomes Dwyer maps as the class of monomorphisms he had been missing: his account describes an open immersion whose map into its closure has a right adjoint, and observes that this implies asphericity. He conjectures that the weaker asphericity condition might itself suffice for preservation of weak equivalences under cobase change, and states the dual proposal for closed immersions and left adjoints. These are proposed extensions, not established results on the page. He hopes they will lead to a more direct model structure on categories and defers the investigation to part~V. The marginal circled interior symbol and the pen repairs in the discussion are retained (p.\ 291).
Thomason's remarks on algebraic $K$-theory lead to Anderson's work and to the author's own “integration” and “cointegration” of homotopy types, projected for part~VI on derivators. Grothendieck wants to examine Anderson's indexing-category assumptions, but argues that the notion of a homotopy limit should depend only on weak equivalence. For the direct homotopy limits arising in the cited application, he points to the explicit categorical construction sketched in \S69, pp.\ 198--199, and regards the model structure as unnecessary for that purpose. Cointegration appears less direct in $(\mathrm{Cat})$; he contrasts this with the anticipated geometric behavior of higher groupoid stacks and expects fibrations, characterized intrinsically by preservation of weak equivalence under base change, to enter its categorical construction. The references and evaluations here are the author's account of the papers, not an independent verification of their claims. Section~88 then pauses the development to inventory the questions still outstanding before part~V (p.\ 292).

\paragraph{\S88. An inventory of unfinished questions (pp.\ 292--300).}
Before beginning part~V, Grothendieck lists five subjects still requiring attention. The first is the direct-image problem of \S84: whether an appropriate $f_!$ preserves aspheric objects. His earlier suspicion of counterexamples has not become a result, and he now allows that an unexpected positive statement may exist. The second is the promised systematic treatment of contractors and generators for contractibility, parallel to aspherators for asphericity. Apart from the trivial contractor equivalent to a final category, their strict-test property connects their presheaf homotopy formalism with pure homotopy theory (pp.\ 292--293). The third asks for an autodual homotopy formalism. Interval homotopy degenerates in a pointed category, whereas complexes in an abelian category supply meaningful, autodual notions of homotopy, homotopism, null-homotopy and quasi-isomorphism. The mapping-cylinder comparison belongs within each of these structures; it does not identify the two kinds of equivalence. The analogy between $H^0$ with its shifts and connected components suggests a broader formalism, possibly involving fibers, cofibers and exact sequences. This programme is linked to the integration and cointegration of homotopy types, the discussion of derivators in \S69, and programme C of \S71, pp.\ 207--210 (pp.\ 293--295). A nearer question concerns the induced structure on $M_{/a}$: the usual aspheric structure on a presheaf slice can fail total asphericity even for a contractor, suggesting that the appropriate equivalences should compare homotopy types fiber by fiber (pp.\ 295--296).

The fourth item becomes an extended digression on why semisimplicial homotopy is preserved by an arbitrary functor $A\to B$, although the interval formalism of \S51 assumed preservation of finite products. When products and distributive sums exist, a simplicial set $I$ supplies objects $I_A$, and maps $I_A\times X\to Y$ can be described degreewise by maps from $I_n$ into the sets of morphisms $X_n\to Y_n$, with the corresponding compatibility. The latter description makes sense without those exactness hypotheses. Fully faithful embeddings preserve it; embedding $A$ and $B$ in presheaf categories explains why the induced $f_!$ preserves the particular products $I_A\times X$, since they are sums of copies of $X$, without preserving every finite product. The author therefore contemplates homotopy and contractibility structures, and their morphisms, without global products or final objects. He compares this enlargement with the earlier passage to sites, but does not regard it as urgent (pp.\ 296--299). The checked text retains the differing visible variance marks in the Hom displays rather than harmonizing them.

The fifth item asks for intrinsic criteria making categories specified by algebraic structure types into canonical modelizers. Small limits and colimits, internal Hom objects in slices, universal and disjoint sums, total 0-connectedness and sections of nonempty objects form the provisional list of properties. The expected examples include elementary presheaf constructions and categories of $n$-stacks and $\mathrm{oo}$-stacks. The desired conclusion is a canonical contractibility structure whose associated usual weak-equivalence structure is modelizing, and then the analogous assertion for every basic localizer. At the end, a separating contractible interval is proposed as a deliberately tentative test conjecture; the author would first seek a counterexample. None of this is presented as an established characterization (pp.\ 299--300). Section~89 interrupts the inventory to repair a specific claim made in \S85 about strict generation.




\paragraph{\S89. Strict generation and the topos obstruction to left exactness (pp.\ 300--302).}
Returning to \S85 after three days, Grothendieck reexamines the claim that a full subcategory of $(\mathrm{Cat})$ containing the standard simplices must be strictly generating. Enlarging a strictly epimorphic family need not preserve strict epimorphicity, so the earlier “a fortiori” was insufficient. He first invokes the stronger universally strict generation furnished by the low-dimensional simplices, citing Giraud's descent article, Proposition~2.5, p.\ 28. He then records a direct argument that enlargement does preserve strict generation for full subcategories, despite the failure for arbitrary families (pp.\ 300--301). The reference to the Glossaire in SGA~4 distinguishes the relevant epimorphism notions. These observations repair the statement on p.\ 284 rather than validate its original justification.
The section then settles the earlier doubt about the left adjoint of $\Delta\hookrightarrow(\mathrm{Cat})$. For a cocomplete category $M$ and a small strictly generating full subcategory $A$, the proposition says that $i_!:A^\wedge\to M$ is left exact if and only if $M$ is a topos. The proof identifies $M$ with the full image of its fully faithful $i^*$ and applies Giraud's criterion to the left adjoint of that inclusion. Since $(\mathrm{Cat})$ is not a topos, its $i_!$ is not left exact; the text points to epimorphisms that are not strict or effective (p.\ 301). The corollary locates the failure in fibered products $b\times_F c$ with $b,c$ in $A$ and general $F$ in $A^\wedge$, even while finite products and fibered products over an object of $A$ are preserved in the situation of \S85. This answers the sharper-version question in Remark~1, p.\ 283: additional hypotheses are required (p.\ 302). The explanation on p.\ 302 literally says “if $M$ is strictly generating”; the edition preserves that wording, alongside the proposition's earlier designation of $A$ as the strictly generating subcategory. Section~90 then resumes the unfinished inventory.




\paragraph{\S90. Existence and perfection of test functors,\newline followed by an afterthought on strict generation (pp.\ 302--312).}
The sixth item in the unfinished inventory asks for test functors from a test category $A$ into a modelizing asphericity structure $M$. The usual finest localizer is the strongest case: a functor satisfying the corresponding test conditions also serves for coarser basic localizers. For contractibility structures, the preferred functor lands in $M_c$ and takes the contractible generators of $M$ to contractible presheaves under $i^*$, thereby supplying the bimorphism setting of \S86. This reformulation separates the existence question from the ambient category as far as possible and suggests testing on a small collection of generators, perhaps a single interval (pp.\ 302--303). The image phrase on p.\ 303 literally says “image of A by A”; the checked text does not silently replace its last letter by the witness's $i$.

The value of a direct test functor is stronger than the abstract availability of modelizing comparisons. A composite through a presheaf category and $(\mathrm{Cat})$ already transfers homotopy types under the earlier hypotheses, but a direct $i:A\to M$ also gives the external object $\mathrm{Hom}^{(A)}(x,y)$ with value $\mathrm{Hom}_M(i(a)\times x,y)$ at $a$. The discussion links this to Quillen's semisimplicial categories and, through the presheaf enlargement of $M$, avoids assuming all internal Hom objects exist in $M$ itself. Normalizing global sections gives the final-object conditions in (5)--(7). The distinction between a left-exact functor into a slice $A^\wedge/E$ and one into $A^\wedge$ matters here; the earlier composite need not have the latter property (pp.\ 304--305).

A further demand is that the left adjoint $i_!$ also modelize homotopy types. The sharper condition that every unit $F\to i^*i_!F$ be a weak equivalence introduces a \emph{perfect test functor}; without the modelizing assumptions, the parallel expression is a perfect aspheric functor. These are intended to give mutually inverse localized adjoints. The ordinary inclusion $\Delta\to(\mathrm{Cat})$ fails this stronger condition, whereas the canonical slice functor associated with a weak test category and Thomason's subdivided construction provide motivating positive examples. The author asks how far these examples extend to arbitrary test categories and contractors, and what special features of $\Delta$ would be responsible if they do not (p.\ 306). This leads to the category $M^A$, its alternative interpretation when $M=B^\wedge$, and the need to distinguish the weak-equivalence structures one can put on it (p.\ 307). The first occurrence of “B is a test functor” is retained literally, while a later occurrence is visibly corrected to “test category.”

The dual interpretation of $M^A$ suggests importing Quillen's projective construction for $M^o$ as an injective construction for $M$. Factoring the map from the initial to the final diagram as a cofibration followed by a trivial fibration is proposed as a way to produce a test functor, with the author's explicit qualification “more or less, at any rate imply.” It is a proposed escape from the stalled inductive construction, not a completed proof. Two further questions ask whether Quillen's theorems extend from $\Delta$ to general test categories and whether the localized diagram categories are canonically independent of that choice. The topos case is compared with Illusie's semisimplicial sheaves and projected for part VII. Grothendieck then notices that the first question depends on an asphericity or contractibility structure while the others appear to depend only on the category, and acknowledges that the questions remain insufficiently clear (pp.\ 308--309).

The dated comments of 4 July reduce existence questions to elementary presheaf modelizers when a suitable testing functor is already available, with a corresponding reduction in the perfect case. For the standard inclusion $i$, the functors $i_!\mathrm{Sd}^n\alpha$ for $n\geq2$ give the Thomason example; $\alpha:\Delta\to\Delta^\wedge$ is the canonical inclusion. The notation changes from the author's earlier Bar to Thomason's Sd. Grothendieck speculates that analogous subdivision operations, perhaps arising from closed-model factorization, could turn ordinary test functors into perfect ones. This places the existence problem firmly in the projected part V rather than resolving it here (pp.\ 309--310).

After the asterism, an unnumbered afterthought returns to \S89. A family $u_i:X_i\to X$ is strictly epimorphic when the induced Hom map is injective and every compatible family, tested against all common-source squares, comes from a map out of $X$. The property depends on the generated sieve, but strictness alone need not survive enlargement. Universal strictness requires stability under arbitrary base change and does survive enlargement; when fibered products exist, the compatibility test reduces to $X_i\times_X X_j$. A direct-limit cocone is strictly epimorphic (pp.\ 310--311). The proposition then equates four conditions for a small full $A\subset M$: strict epimorphicity of every family of arrows from $A$, universal strict epimorphicity of those families, their presentation of each object as the indicated direct limit, and full faithfulness of $i^*$. This supplies the definition of a strictly generating subcategory and proves the enlargement corollary by the second criterion. The text cites SGA IV, expos\'e I, Proposition 7.2, p.\ 47, for part of the equivalence and leaves the proof to the reader. Its literal “other three criteria” and the sentence ending the digression are retained. Section~91 begins immediately below (p.\ 312).



\paragraph{\S91. Finite type, boundary operations, and the remaining questions from part II (pp.\ 312--319).}
The seventh item compares finite triangulations of spaces with objects of finite presentation in $\Delta^\wedge$. Grothendieck sketches why the two descriptions should determine the same homotopy types: an ordering of vertices gives one direction, while induction on simplices, piecewise-linear realization and perhaps two barycentric subdivisions suggest the other. The pen additions explicitly require a compact piecewise-linear structure and finite triangulation. For a general category admitting filtering direct limits, finite presentation means that the covariant Hom functor commutes with those limits. In $A^\wedge$, the finite-presentation objects are finite direct limits of representables; the author describes this as the 2-universal problem of adjoining finite direct limits to $A$. The question is when these objects describe exactly the finite-type homotopy types (pp.\ 312--313).

The abandoned opening of p.\ 314 is replaced by a proposed counterexample. A finite group $G$ acts freely on an aspheric space $E_G$, with quotient $B_G$. A suitably chosen small full category of aspheric spaces, containing the interval and $E_G$ and closed under finite products, supplies the test category. The finitely presented quotient presheaf $F=E_G/G$ is identified with $i^*(B_G)$; the author argues that its homotopy type is not of finite type when $G\ne1$, using nonvanishing cohomology in arbitrarily high dimensions. This demonstrates why extra restrictions on $A$ are needed in the proposed correspondence. He suggests that quotients of representables remain representable, that each representable have only finitely many sieves, and perhaps that Hom sets be finite. For the converse inclusion, a perfect test functor $\Delta\to A^\wedge$ with finitely presented values would suffice, linking this question back to \S90. These conditions are proposals rather than a proved characterization. The author expects the familiar simplicial, cubical and hemispherical categories, and their finite products, to satisfy the desired finiteness property. He again resists treating the simplicial category as uniquely privileged, while allowing the strict Dold--Puppe theorem as a possible special feature of the familiar three (pp.\ 314--315).

The eighth item asks for boundary operations intrinsic to a test category. Computing through a test functor from a familiar category, or through a nerve, does not answer that question. The proposed geometric data are a functor $i:A\to(\mathrm{Ord})$ with finite values, each having a largest element and realizing as a cell, with the complement of that largest element realizing as its boundary sphere. Principal lower ideals must themselves be represented by objects of $A$, through maps $b\to a$; uniqueness in the slice gives the face objects $a_x$ and maps $\partial_x$ for codimension-one elements. Orientations $\omega_a$ then define signs $\varepsilon_a(x)$ by the induced boundary orientation “à la Stokes” (pp.\ 315--317). With the dimension decomposition of $A$, a contravariant additive-valued functor produces the displayed sums $K_n=\coprod_{a\in A_n}K.(a)$ and their signed boundary maps. The manuscript still asks for natural extra conditions ensuring that these complexes calculate the ordinary homology and cohomology of the associated category, also with twisted coefficients, and ensuring that $A$ is a test category. The formulas therefore describe a tentative construction, not a completed comparison theorem (pp.\ 316--317).

Degeneracy operations are the next obstacle. The first attempted example is boxed and crossed out on p.\ 317; the composition display inside it has genuinely blank arrow positions. On p.\ 318 Grothendieck explains the failure. He starts with a full category $A_o$ of finite ordered sets satisfying the cell condition, stable under products and passage to principal ideals, and containing the three-point ordered set whose two outer points map to the middle point. He then tries objects consisting of such an ordered set and an orientation. Allowing all underlying maps would preserve the desired product argument but fail invariance of orientations under isomorphisms, an assumption now recognized as necessary for the differential. Enforcing that compatibility destroys the available product and representability argument for the interval. Thus the intended test-functor proof does not go through. The difficulty leads to the opposite possibility: boundary operations might intrinsically single out the simplicial, cubical and hemispherical examples, just as the Dold--Puppe question had suggested earlier. This possibility remains open (pp.\ 317--318).

The ninth item returns to \S43. The category of simplices without degeneracies is asserted to be a weak test category but not a test category; the author wants to write down the proof and consider the analogous cubical and hemispherical assertions. He also notes that the hemispherical category itself has not yet been formally introduced, and that its strict-test or stronger contractor property still deserves proof. Six comparison test functors between the three familiar types of complexes would be useful, but he may omit them from a work that is not intended to be an “Elements d'Algèbre Homotopique” (pp.\ 318--319). The other recalled question concerns the ordered set of nonempty finite subsets of an infinite set: a positive weak-test result would support the modelizing role of $(\mathrm{Ord})$. He asks whether that modelizing structure comes from an asphericity structure or even a contractibility structure, and poses the same issue in the general situation of the earlier proposition on p.\ 74. Finally, for basic localizers $\underline W\subset\underline W'$ satisfying Loc~4, he asks whether total $\underline W$-asphericity and total $\underline W'$-asphericity are equivalent, or equivalently whether either is the same as total 0-connectedness. These are the concluding questions of the cote, not claims established by it (p.\ 319).

